\documentclass[../main.tex]{subfiles}
\begin{document}
%==========================================TIÊU ĐỀ TẬP========================================
\begin{table}[h]
    \begin{adjustwidth}{-1cm}{}
        \begin{tabular}{>{\centering\arraybackslash}p{6cm}|>{\raggedright\arraybackslash}p{12.5cm}}
            \multirow{5}{6cm}
            % Cột bên trái: ÉPISODE và số tập
            {\\[-40pt]
                \flaregothic\fontsize{20pt}{20pt}\selectfont \centering
                \color{eptype}ÉPISODE\color{black}\\
                \burbank\fontsize{72pt}{72pt}\selectfont
                \color{epnum}14\color{black}\\[6pt]
                \cafeteria\fontsize{20pt}{20pt}\selectfont\color{epnum}(1-10)\color{black}
                \\[18pt]
                \flaregothic\fontsize{18pt}{18pt}\selectfont
                \color{epattr}ÉPISODE PERDU
                \color{black}
            }
             &
            % Dòng thứ nhất: tên tập tiếng Nhật
            {\fotskip\fontsize{18pt}{18pt}\selectfont サメハダー vs. アイドル
            }                                  \\[6pt]
            % Dòng thứ hai: tên tập tiếng Anh
             & {\cafeteria\fontsize{18pt}{18pt}\selectfont Sharpedo vs. \hll}                        \\[4pt]
            % Dòng thứ ba: tên tập tiếng Việt
             & {\vnmsans\fontsize{18pt}{18pt}\selectfont Bãi biển trong nhà}                         \\[8pt]
            % Dòng thứ tư: Nhân vật xuất hiện
             & {Nhân vật: \emp{kuon1} \emp{miko} \emp{aki} \emp{matsuri} \emp{haato} \emp{ayame}}
            \\[8pt]
            % Dòng cuối cùng: Thời gian diễn ra của tập (thường sẽ là ngày phát hành)
             & {\continuum\fontsize{16pt}{16pt}\selectfont Ngày phát hành: 11/08/2019}               \\[18pt]
        \end{tabular}
    \end{adjustwidth}
\end{table}
%=============
\txtbx[2]{{\color{vol} \uline{Tóm tắt:}} Giữa cái nóng oi bức của mùa hè, Ayame nảy ra sáng kiến tự dựng một bãi biển siêu mini ngay giữa văn phòng để mọi người cùng diện đồ bơi mát mẻ. Rắc rối ập đến khi một con cá mập Sharpedo bất ngờ trồi lên từ hồ bơi, khơi mào cho cuộc chiến Cá xoáy làm đổ sụp cả cửa kính văn phòng. Sau sự kiện hỗn loạn, một tin vui bất ngờ tới khi Kuon nhận được cuộc gọi từ gia đình rủ đi biển thực tế, mở đường cho chuyến đi nghỉ hè đầy ắp tiếng cười của đại gia đình \hll.}

%=========================================TRUYỆN CHÍNH========================================
\color{black}

Một buổi trưa hè tháng Tám, cái nắng gay gắt như thiêu như đốt khiến ai nấy đều cảm thấy bức bối và khó chịu.

Văn phòng \hll\ hôm nay bỗng trở nên ồn ào và náo nhiệt hơn mọi ngày, nhưng chẳng phải vì sự bận rộn của công việc, mà là bởi cái nóng nực oi bức bủa vây. Hệ thống điều hòa nhiệt độ vẫn bặt vô âm tín kể từ khi nơi này khai trương vào cuối tháng Năm, và tất cả những gì mọi người có được chỉ là vài chiếc quạt máy cũ kỹ chạy ì ạch, hòa lẫn trong tiếng quạt giấy phành phạch đầy mệt mỏi của các thành viên.

\begin{figure}[H]
    \centering
    \begin{subfigure}[b]{0.45\textwidth}
        \centering
        \includegraphics[width=8cm]{../../../../Image/Book1/1-10a1.png}
    \end{subfigure}
    \quad
    \begin{subfigure}[b]{0.45\textwidth}
        \centering
        \includegraphics[width=8cm]{../../../../Image/Book1/1-10a2.png}
    \end{subfigure}
\end{figure}

Matsuri, diện một bộ đồ bơi hai mảnh màu cam rực rỡ, đứng tựa người vào khung cửa sổ, ánh mắt mơ màng ngắm nhìn bầu trời xanh thẳm ngoài kia. Mái tóc nâu sáng của cô lấp lánh dưới ánh dương, nhưng nét mặt lại đượm một nỗi buồn chán chường:
``Mình muốn đi biển quá đi mất\dots''

Aki và Miko đang ngồi thong thả trên chiếc ghế bành gần đó, tay không ngừng phe phẩy chiếc quạt giấy xua đi cái nóng, uể oải đồng thanh phụ họa:
``Nhất trí\dots''

Kuon, vừa vươn vai đặt bút xuống sau khi hoàn tất bản kế hoạch chuẩn bị cho dự án \hll\ Summer, ngả người tựa lưng vào ghế. Cậu cầm cốc nước đá lạnh buốt uống ừng ực một ngụm lớn rồi bộc bạch đầy tiếc nuối:
``Anh cũng muốn đi giải nhiệt lắm chứ, nhưng mà lấy đâu ra ngân sách bây giờ? Doanh thu từ đợt bán hàng lưu niệm ở Comiket còn chưa được chuyển về tài khoản công ty nữa là\dots''

Nghe vậy, nàng Cổ động liền ba chân bốn cẳng chạy ào tới trước mặt ba người. Cô dang rộng hai tay ra để minh họa cho sự bức bách tột độ của mình:
``Mọi người nhìn em nè! Bộ đồ bơi mới tinh của em đang khao khát được diện ra biển đến mức sắp khóc thét lên rồi đây này!''

Từ góc phòng, Haato - người nãy giờ vẫn đang cắm cúi làm việc riêng để giết thời gian - ngước lên nhìn Matsuri bằng ánh mắt tò mò không giấu giếm:
``Vậy ra đó là lý do bà chị tự nhiên lại mặc đồ bơi đi vòng vòng quanh văn phòng từ sáng đến giờ à?''

Cậu Tổng liếc nhìn cô nàng đầy hoài nghi:
``Sao nghe lý do đó nó vô lý đến thế nhỉ\dots\ Nhưng mà mặc thử là một chuyện, còn lôi nhau ra bãi biển thật lại là chuyện hoàn toàn khác đấy,'' cậu khẽ lắc đầu, tặc lưỡi cảm thán: ``Vụ này xem chừng khoai thật à nha\dots''

Dường như với tình hình tài chính eo hẹp hiện tại của công ty, việc tổ chức một chuyến đi nghỉ mát ở biển là điều hoàn toàn bất khả thi. Cả văn phòng bỗng chìm vào khoảng lặng. Thế nhưng, ngay lúc đó, một giọng nói trong trẻo bất chợt vang lên, phá vỡ bầu không khí tĩnh mịch:
``Nếu vậy thì\dots\ sao chúng ta không tự 'dựng' luôn một bãi biển ở đây nhỉ?''

Nàng Quỷ Ayame đột ngột xuất hiện giữa phòng, híp mắt nở một nụ cười rạng rỡ đầy hứng khởi. Matsuri và Aki đồng loạt trố mắt ngạc nhiên, trong khi Kuon chỉ biết nheo mắt nhìn cô nàng bằng ánh mắt nghi hoặc tột độ:
``Bằng cách nào cơ? Đừng bảo là em định hì hục xúc cát với bơm nước biển từ ngoài vịnh chở về đây đấy nhé? Cái vụ tràn dầu làm ngập cả văn phòng tháng trước vẫn còn sờ sờ ra đấy, anh chưa quên đâu!''

Nàng Oni che miệng cười khúc khích, đáp lại bằng một giọng điệu đầy tự tin và kiêu hãnh:
``Tất nhiên là không phải trò thủ công ngốc nghếch đó rồi! Mọi người cứ ngồi đó mà chờ xem phép màu của em nhé!''

Dứt lời, cô nàng ba chân bốn cẳng chạy tót ra ngoài hành lang, dáng vẻ vội vã và hào hứng như thể đang ấp ủ một kế hoạch vô cùng vĩ đại. Cậu Tổng đứng đó, đưa tay gãi đầu khó hiểu, thầm cầu nguyện trong bụng: ``\textit{Chỉ mong là Ojou-san đừng làm điều gì điên rồ quá đáng để anh phải nai lưng ra dọn dẹp hậu quả\dots}''

\ct

Chẳng bao lâu sau, nàng Quỷ đã hớn hở quay trở lại văn phòng cùng với thành quả vĩ đại của mình: một ``bãi biển siêu mini'' được thiết kế cực kỳ đơn sơ từ một chiếc phao bơi nhỏ xíu bơm đầy nước, được trang trí xung quanh bởi những tấm bìa các-tông in hình bãi biển thật thụ mà cô nàng vừa hì hục cắt ghép.

\img{1-10b}

``Xong xuôi hoàn hảo rồi đây!'' Ayame reo lên tự hào. 

Lúc này, các thành viên khác cũng đã hào hứng khoác lên mình những bộ đồ bơi lộng lẫy đủ màu sắc. Nàng Dị giới Aki vỗ tay tán thưởng, ánh mắt không giấu nổi sự kinh ngạc:
``Trời ạ, em làm nhanh tay thật đấy Ayame-chan!''

Dù vậy, cậu Tổng vẫn đứng khoanh tay, đôi mắt dò xét đánh giá tác phẩm nghệ thuật kỳ quặc trước mặt:
``Ờm\dots\ Thế này mà cũng được coi là đi biển luôn à\dots?''

Nàng Quỷ bật cười sảng khoái, chống hai tay lên hông, tự tin đáp trả với việc nhấn mạnh từng từ ngữ:
``Thì ban nãy Matsuri-chan chỉ khóc lóc bảo là muốn \uline{đi biển} thôi mà, chứ cậu ấy có nói rõ là phải \uline{tắm biển} đâu cơ chứ.'' 

Nghe lý luận cùn nhưng vô cùng thuyết phục ấy, Kuon chỉ biết gật gù bó tay.

Trái ngược với sự hào hứng của mọi người, nàng Cổ động ngồi xổm trước ``bãi biển'', ủ rũ cúi gầm mặt xuống, thở dài thất vọng:
``Nhìn cái đống này chả giống chút nào so với những gì tôi tưởng tượng và kỳ vọng cả\dots''

Aki bước tới đứng cạnh Ayame, vẻ mặt đầy mãn nguyện và tươi rói, dịu dàng động viên cô bạn:
``Thôi nào, thế này cũng được tính là có không khí rồi. Với cả cũng lâu lắm rồi tôi mới có dịp được mặc đồ bơi vui vẻ cùng mọi người thế này mà!''

\begin{figure}[H]
    \centering
    \begin{subfigure}[b]{0.45\textwidth}
        \centering
        \includegraphics[width=8cm]{../../../../Image/Book1/1-10c1.png}
    \end{subfigure}
    \quad
    \begin{subfigure}[b]{0.45\textwidth}
        \centering
        \includegraphics[width=8cm]{../../../../Image/Book1/1-10c2.png}
    \end{subfigure}
    \quad
    \begin{subfigure}[b]{0.45\textwidth}
        \centering
        \includegraphics[width=8cm]{../../../../Image/Book1/1-10c3.png}
    \end{subfigure}
    \caption{Các thành viên trong trang phục áo tắm.}
\end{figure}

Cậu Tổng vừa lau những giọt mồ hôi lấm tấm trên trán, vừa gật gù đồng ý:
``Ừ thì, ít nhất thì cái trò này cũng giúp mọi người đổi gió thư giãn được một chút. Cơ mà\dots'' 

Vừa nói dứt lời, cậu vội vàng đưa tay lên che nửa khuôn mặt đang đỏ ửng lên như trái gấc chín của mình. Tinh mắt phát hiện ra biểu cảm ngượng ngùng ấy, nàng Oni lập tức sấn tới, chống nạnh và nở một nụ cười tinh quái trêu ghẹo:
``Ái chà\~{}? Sao tự nhiên mặt anh lại đỏ rực lên như tôm luộc thế kia hả Kuon-senpai\~{}?''

Cậu Tổng bối rối lúng túng chưa kịp mở miệng thanh minh thì nàng Cổ động Matsuri bỗng nhiên đứng phắt dậy, lớn tiếng tố cáo đầy phấn khích:
``Là vì anh ấy rất thích được ngắm nhìn tụi mình mặc đồ bơi đó! Mới hôm trước ở sự kiện anh ấy vừa tự nhậnnhận--''

Nhanh như một cơn gió lốc, Kuon lập tức lao tới dùng tay bịt chặt lấy miệng cô nàng lém lỉnh, cố gắng dập tắt ngọn lửa hóng hớt trước khi nó kịp bùng lên thiêu rụi hình tượng nghiêm túc của cậu. Mặt cậu đỏ bừng bừng, hốt hoảng nạt nhỏ:
``Trời đất ơi, sao em lại ăn nói lớn tiếng bô bô ra như thế hả!?''

Mặc cho miệng đang bị bịt kín mít, đôi mắt Matsuri vẫn ánh lên nét tinh nghịch, cô cố gắng ú ớ trêu chọc:
``Ứ\dots\ Ì\dots\ ữa\dots\ ớ\dots\ anh ả\dots\ ự\dots\ ừa\dots\ ận\dots\ ấy\dots\ ôi\dots'' (Thì bữa trước anh chả tự thừa nhận đấy thôi.)

Cậu Tổng nghiến răng đe dọa, cố gắng duy trì chút uy nghiêm cuối cùng của một quản lý:
``Đừng có nói thêm chữ nào nữa, nghe rõ chưa!? Cấm có hé nửa lời! Nếu không là anh trừ nửa tháng lương của em đấy!'' 

Dưới sự đe dọa kinh tế đầy uy lực đó, nàng Cổ động đành phải ngoan ngoãn gật đầu đồng ý. Cậu thở phào nhẹ nhõm buông tay ra. Thế nhưng, ngay giây phút ấy, cậu bỗng cảm nhận được những ánh mắt sắc như dao cạo từ các cô gái trong phòng đang đổ dồn về phía mình với đủ mọi sắc thái biểu cảm:

\begin{itemize}
    \item Aki cười mỉm nhẹ nhàng, ánh mắt vô cùng thích thú như đang xem kịch hay.
    \item Ayame thì nhìn cậu với vẻ mặt vô cùng đắc ý, như thể vừa tóm được cái thóp chí mạng.
    \item Haato tuy ngồi xa nhưng cũng lén lấy tay che miệng cười khúc khích hóng chuyện.
\end{itemize}

Bất lực trước tình thế muối mặt, Kuon chỉ biết đưa tay ôm trán, thở dài thườn thượt:
``Tuyệt vời thật đấy. Em vừa mới đào hố chôn sống anh luôn rồi đó Matsuri ạ\dots''

\ct

Một lát sau, Miko bước vào văn phòng với bộ đồ bơi dễ thương đã được diện sẵn. Nàng Vu nữ hớn hở chuẩn bị tận hưởng buổi tắm hồ bơi mini do Ayame thiết kế:
``Được rồi mọi người ơi, tất cả đã chuẩn bị sẵn sàng chưa nhenhe--''

\img{1-10d}

Thế nhưng, khi vừa nhìn thấy các cô gái ai nấy đều diện những bộ áo tắm quyến rũ và lộng lẫy, Miko bỗng ôm chặt lấy ngực trái, chân nam đá chân chiêu rồi khuỵu hẳn xuống sàn với vẻ mặt vô cùng mãn nguyện. Cô lầm bầm trong sự xúc động nghẹn ngào:
``Đồ bơi của mọi người\dots\ sao mà đáng yêu và xinh xắn quá trời quá đất luôn vậy\dots'' 

Cậu Tổng hoảng hốt vội vàng chạy tới, nắm lấy cánh tay cô đỡ dậy, dỗ dành:
``Thôi nào, anh thấy bộ đồ của em cũng rất đẹp và hợp mà. Nào, mau đứng lên đi em.''

Thấy mọi người đã tề tựu đông đủ quanh chiếc phao bơi nhỏ, Ayame hắng giọng, nở một nụ cười rạng rỡ đầy khí thế:
``Giờ thì tất cả cùng quẩy lên thôithôi--'' 

Nhưng còn chưa kịp dứt lời, cô nàng bỗng khựng lại. Ánh mắt cô dán chặt vào một vật thể kỳ lạ màu xám đang từ từ nhô lên từ giữa hồ bơi mini. Giọng Ayame run run đầy bối rối:
``Ể\dots\ C-Cái gì thế kia!?''

\img{1-10e}

Kuon nheo mắt nhìn kỹ vật thể nhọn hoắt ấy, rồi bất giác giật mình lùi lại vài bước cảnh giác:
``Anh nhìn không nhầm thì đó hình như là vây lưng của một con cá mập\dots\ Nhưng mà làm quái nào nó lại chui được vào cái hồ bơi phao bé xíu này chứ?'' 

Miko, sau khi đã lấy lại được bình tĩnh, vuốt cằm và cất giọng giải thích bằng một tông giọng sặc mùi ``uyên thâm'' học thuật:
``Chị đã từng xem qua cảnh này rồi. Theo nguyên lý điện ảnh kinh điển, trong một phân cảnh nếu hội tụ đủ ba yếu tố: bãi cát, nước biển, và các mỹ nhân diện đồ bơi bốc lửa cùng xuất hiện tại một thời điểm\dots\ thì chắc chắn cá mập sẽ bị thu hút và được triệu hồi đến.''

\begin{figure}[H]
    \centering
    \begin{subfigure}[b]{0.45\textwidth}
        \centering
        \includegraphics[width=8cm]{../../../../Image/Book1/1-10f1.png}
    \end{subfigure}
    \quad
    \begin{subfigure}[b]{0.45\textwidth}
        \centering
        \includegraphics[width=8cm]{../../../../Image/Book1/1-10f2.png}
    \end{subfigure}
    \caption{Miko `ưu tú' đưa ra giả thuyết về con cá mập xuất hiện trong hồ bơi, dựa trên kinh nghiệm xem phim của cô.}
\end{figure}

Nàng Vu nữ chỉ tay thẳng về phía chiếc vây cá mập đang nhấp nhô giữa dòng nước cạn, khẳng định chắc nịch:
``Nói tóm lại, đó đích thị là Cá mập \hll!'' 

Nghe xong lời phán quyết xanh rờn ấy, Matsuri, Aki và Haato đồng loạt hét lên đầy kinh hãi, hoảng loạn lặp lại từ ``Cá mập''. Trái ngược với phản ứng của họ, Kuon chỉ biết đưa tay xoa thái dương, ngán ngẩm thở dài:
``Anh nghĩ là em dạo này bị nhiễm dăm ba cái kịch bản phim hạng B hơi bị nặng rồi đấy Miko-chi ạ\dots''

Miko vẫn hất cằm quả quyết:
``Em có lý luận khoa học hẳn hoi nhé! Anh không tin thì thôi, mặc xác anh!'' 

Trong khi mọi người còn đang tranh cãi, nàng Quỷ Ayame - với bản tính hiếu kỳ không sợ trời không sợ đất - lại từ từ tiến sát lại gần hồ bơi. Cô nàng nghiêng đầu, tò mò hỏi một câu vô thưởng vô phạt:
``Nhưng mà chẳng phải cái tên gọi `Cá mập Công sở' nghe nó ngầu đét và kêu hơn hẳn sao?''

Đúng lúc đó, con cá mập đột ngột trồi hẳn lên khỏi mặt nước. Nó há to cái miệng đầy những chiếc răng cưa sắc nhọn và\dots\ \textit{PHẬP!} Ngoạm thẳng vào bắp chân của nàng Quỷ. Con quái thú dùng sức kéo mạnh cô xuống nước, khiến cô nàng hoảng hốt la oai oái, kéo theo cả căn phòng rơi vào cảnh náo loạn tột độ.

\img{1-10g}

``\textbf{Ayame / Ayame-chan! / Ojou-san!!}'' tất cả mọi người đồng thanh hét lên kinh hoàng.

Thế nhưng, có một điều vô cùng kỳ lạ là nàng Quỷ không hề cảm thấy chút đau đớn nào. Khi ngoái đầu liếc nhìn xuống chân mình, cô mới nhận ra sự bất thường.
``Ủa? Hình như nó không thực sự cắn rách da mình này!'' 

Bất ngờ thay, con cá mập bỗng cất tiếng khóc thút thít. Giọng nó yếu ớt và tội nghiệp như đang nài nỉ:
``Tôi đang bị trật khớp hàm đau lắm mấy người có biết không\dots''

Matsuri cúi người xuống, khuôn mặt hiện rõ vẻ đồng cảm thương xót:
``Trời đất, ngươi có sao không vậy cá mập?'' 

Con cá mập ngước đôi mắt ươn ướt lên, đuôi vẫy vẫy nhẹ trên sàn, mếu máo trả lời:
``Tại vì dạo gần đây tôi toàn phải ăn sứa biển độc thay cơm ấy mà\dots'' 

Cậu Tổng khoanh tay tặc lưỡi, nhướng mày châm chọc một câu phũ phàng:
``\dots Nên cái mỏ của ngươi bị nọc độc làm cho tê liệt hết cả hàm rồi chứ gì? Mà khoan đã, ta tưởng cá mập thì phải đi săn cá lớn mà ăn chứ, đi ăn sứa làm cái quái gì?''

Con cá mập không thèm đáp lời, chỉ biết giữ im lặng ngậm ngùi vẫy đuôi. Trong khi đó, Ayame - người lẽ ra phải hoảng sợ vì bị cá mập cắn - lại bất ngờ nằm dài ra sàn nhà. Cô nàng tạo dáng thướt tha, khoanh hai tay lại trước ngực rồi nở một nụ cười cực kỳ tự mãn:
``Nè mọi người nhìn thử xem, trông tôi bây giờ có giống nàng tiên cá xinh đẹp trong truyện cổ tích không?''

\img{1-10h}

Kuon toát mồ hôi hột trước sự vô tư đến mức ngốc nghếch của cô nàng. Cậu không ngần ngại vả ngay một gáo nước lạnh:
``Trông chả giống một tẹo nào cả. Nàng tiên cá người ta đuôi thon thả, còn cái 'đuôi cá mập' của em thì nó to sụ béo múp thế kia cơ mà.'' 

Cậu quay sang nhìn quanh mọi người cầu cứu:
``Có ai biết cách nào lôi con cá mập khổng lồ này ra khỏi chân Ayame-chan không?''

Cả phòng lập tức lao xao bàn tán tìm cách giải cứu. Bất chợt, Matsuri cảm thấy một luồng sát khí lạnh lẽo tỏa ra từ phía sau gáy. Cô từ từ quay đầu lại, rồi kinh hãi hét lớn lên với âm lượng cực đại:
``\textbf{CÓ MỘT CON CÁ MẬP KHÁC ĐANG Ở KIA KÌA!!}''

Một con cá mập khổng lồ thứ hai bỗng nhiên trồi lên từ hư không ngay giữa hồ bơi mini. Nó quẫy đuôi nhảy phốc lên bờ và hung hãn lao thẳng về phía đám đông. Cả văn phòng lại thêm một phen nhốn nháo loạn cào cào. Tất cả mọi người chạy toán loạn thét lên vì hoảng sợ, kể cả nàng ``tiên cá'' bất đắc dĩ Ayame cũng cố gắng lết cái chân đang bị cá mập cắn để bỏ trốn.

Mọi thứ diễn ra quá chớp nhoáng, nhưng trong khoảnh khắc quay chậm kịch tính như phim hành động Hollywood:

\gif{1-10j1}{1}{218}

\dots chiếc đuôi của con cá mập đang ngậm chân Ayame đột ngột quật mạnh một cú trời giáng vào mặt con cá mập thứ hai. Bị ăn đau, con thứ hai lập tức há mõm ngoạm chặt lấy đuôi của con thứ nhất. Cú ngoạm đau điếng khiến con đầu tiên phải há mồm nhả chân nàng Quỷ ra, và quay ngoắt lại cắn trả vào đuôi của kẻ thù.

\gif{1-10j2}{1}{94}

Thế là hai con quái thú biển khơi xoắn chặt lấy nhau thành một vòng tròn. Chúng xoay tít mù như một cơn lốc xoáy hung tợn, rồi bất ngờ văng mạnh ra khỏi cửa sổ văn phòng với một tiếng *XOẢNG* chói tai, làm vỡ nát toàn bộ hệ thống cửa kính. 

Ngay khoảnh khắc chúng bay ra ngoài không trung, một dòng chữ lớn màu cam chói lọi hiện lên rực rỡ giữa nền trời:

\begin{center}
    \textbf{{\color{orange}**CÁ XOÁY\\(シャークネード)**}}
\end{center}

%==========================================TRUYỆN PHỤ=========================================
\omk

Mọi người sững sờ đứng chôn chân tại chỗ, ngơ ngác nhìn khung cảnh hoang tàn và ô cửa sổ trống hoác trước mắt. Cậu Tổng là người đầu tiên lấy lại được bình tĩnh. Cậu vội vàng chạy tới bên cạnh Ayame, người vừa mới may mắn thoát khỏi hàm cá mập, ân cần hỏi han:
``Em không sao chứ Ojou-san? Có bị xước xát ở đâu không?'' 

Ayame phủi bụi trên người, gật đầu mỉm cười thản nhiên như chưa từng có chuyện gì xảy ra:
``Ưm. Em vẫn hoàn toàn khỏe mạnh và ổn thỏa nha.''

Cả phòng đồng loạt thở phào nhẹ nhõm. Kuon tuy vẫn còn bực mình vì rắc rối cô nàng gây ra, nhưng cũng đành lắc đầu cười trừ bất lực:
``Anh xin em đấy, lần sau bớt bày trò nghịch ngợm lại giùm anh một cái. Văn phòng công ty làm việc chứ có phải là công viên giải trí thu nhỏ đâu cơ chứ.'' 

Thay vì tỏ ra hối lỗi, nàng Oni lại vui vẻ nhún vai, vênh mặt tự mãn:
``Nhưng rõ ràng là mọi người đã có một tràng cười rất vui vẻ mà, đúng không nào\~{}?''

Sự lạc quan tếu của cô nàng khiến cả phòng bật cười sảng khoái, bầu không khí nặng nề bỗng chốc tan biến. Tuy nhiên, cậu Tổng thì lại méo mặt, không thể nào nhếch mép cười nổi. Cậu tiến lại gần, cúi xuống nhặt vài mảnh kính vỡ trên sàn, xoa trán lầm bầm đầy tuyệt vọng:
``Thôi xong đời anh rồi\dots\ A-chan mà về thấy cảnh tượng hoang tàn này là chị ấy sẽ băm vằm cả lũ ra mất\dots\ Đây đã là ô cửa kính thứ ba bị vỡ tan tành trong tháng này rồi đấy\dots''

Sự lo lắng của quản lý khiến bầu không khí lập tức chùng xuống. Miko khép hờ chiếc quạt giấy trên tay, nghiêng đầu hỏi với giọng điệu e dè:
``Vậy\dots\ nếu bị báo cáo lên trên thì chúng ta sẽ bị xử phạt như thế nào hả anh?''

Kuon nhắm chặt hai mắt, khoanh tay trước ngực suy nghĩ một hồi lâu. Khi mở mắt ra, giọng cậu trầm buồn và đầy tính toán:
``\dots Nếu chiếu theo đúng nội quy quy định của công ty, sẽ có hai phương án giải quyết: hoặc là Ayame sẽ phải tự bỏ tiền túi ra đền bù toàn bộ thiệt hại cho cửa kính, hoặc là mỗi người trong chúng ta, kể cả anan, sẽ phải trích một phần tiền lương tháng này để cùng đền. Nhưng mà\dots''

Cậu ngừng lại, thở dài thườn thượt, vò đầu bứt tai:
``\dots anh thấy cả hai cách này đều không ổn chút nào. Mặc dù Ojou-san là người bày ra cái trò dựng 'bãi biển mini' này\dots'' Cậu chỉ tay về phía chiếc hồ phao giờ đã xì hơi, nằm bẹp dúm dưới sàn: ``\dots nhưng lỗi làm vỡ kính thì đâu hoàn toàn do em ấy, là do hai con cá mập kia tự đánh nhau vỡ kính đấy chứ. Còn nếu bắt tất cả mọi người cùng chịu phạt chung thì lại quá bất công. Cho nên là--''

Kuon chưa kịp nghĩ ra giải pháp thì Matsuri đã vui vẻ vẫy tay cắt ngang:
``Anh đừng lo lắng quá Kuon-senpai! Cái cửa kính đó thực ra mua thay thế rẻ bèo à!''

Cậu Tổng ngẩng phắt đầu lên, trố mắt kinh ngạc:
``Thật á? Em chắc không?'' 

Matsuri gật đầu chắc nịch, cười tươi roi rói:
``Chắc chắn 100\% luôn! Hôm khai trương dự án, lúc Aqua-chan vô tình húc vỡ cửa kính văn phòng, chính em với Roboco-chan đã tự đi mua tấm kính mới về thay mà, nên em nắm rõ giá cả lắm!''

Kuon gật gù, dường như đã trút được một phần gánh nặng ngàn cân. Aki bước lên, nhẹ nhàng tiếp lời góp ý:
``Với lại, em nghĩ trong chuyện này ai cũng có một phần trách nhiệm cả. Dù sao thì chính chúng ta cũng đã hùa theo trò vui của Ayame-chan mà.''

Haato cũng mạnh mẽ gật đầu đồng tình:
``Chính xác! Dù sao thì tự sâu thẳm trong lòng, tất cả chúng ta đều khao khát có một chuyến đi biển giải nhiệt cơ mà.'' 

Thế là sau một hồi bàn bạc rôm rả, cả nhóm đi đến quyết định cuối cùng: cả sáu người sẽ cùng nhau góp tiền để mua cửa kính mới đền cho công ty. Cậu Tổng khoanh hai tay, gật gù mỉm cười hài lòng trước tinh thần đoàn kết của các thành viên:
``Được rồi, thế thì mọi chuyện coi như tạm ổn định. Vậy thì bây giờ--''

\textit{Reng reng reng\dots}

Tiếng chuông điện thoại chói tai từ túi quần bỗng cắt ngang lời cậu. Kuon rút máy ra xem, khẽ nhíu mày:
``Xin lỗi mọi người một lát nhé, mẹ anh gọi điện từ quê lên.'' Cậu quay lưng đi ra góc phòng, bắt máy và trò chuyện bằng tiếng mẹ đẻ.

``\vie{Vâng, con nghe đây thưa mẹ?}''

Vừa áp điện thoại vào tai nghe được vài câu, sắc mặt của cậu con trai lập tức biến đổi chóng mặt, từ vẻ bình thản chuyển sang sửng sốt cực độ.

``\vie{Dạ!? Cả gia đình mình cùng với bên tổng công ty của mẹ sẽ tổ chức thuê nguyên chuyến xe đi nghỉ mát ở bãi biển á?}'' Cậu trợn tròn mắt, giọng nói đầy phấn khích: ``\vie{Ngày kia khởi hành luôn ạ? Vâng\dots\ Dạ?}''

Đột nhiên, Kuon hét toáng lên đầy kích động khiến cả văn phòng giật bắn mình:
``\vie{\textbf{Còn dư hẳn 10 chỗ trống luôn ạ!?} Dạ vâng, tuyệt vời quá, để con rủ thêm mấy đứa bạn của con-- Không, không phải bạn gái đâu mẹ! Chỉ là mấy đứa đồng nghiệp làm chung công ty thôi! Vâng, vâng, thế nhé mẹ! Con cảm ơn mẹ nhiều lắm!}''

Cậu cúp máy cái rụp, quay ngoắt người lại đối diện với các cô gái, nở một nụ cười rạng rỡ như ánh mặt trời. Cả đám vẫn đang xì xào tò mò không hiểu quản lý vừa nói chuyện gì mà lại có vẻ vui sướng tột độ đến vậy.

Haato lên tiếng thắc mắc:
``Kuon-senpai, có chuyện gì mà trông anh vui vẻ hạnh phúc thế?'' 

Không thèm úp mở thêm, cậu quản lý dang rộng hai tay lên cao, dõng dạc tuyên bố:
``Các cô gái à, chuẩn bị thu xếp đồ đạc đi! \textbf{Chúng ta sắp được đi biển thật rồi!}''

Cả văn phòng như vỡ òa, mọi người ôm chầm lấy nhau nhảy cẫng lên reo hò trong niềm hân hoan tột độ. Ngay lập tức, kế hoạch cho chuyến đi nghỉ hè hoành tráng được vạch ra chi tiết, đồng thời mọi người cũng gấp rút sắp xếp lại lịch trình để trình đơn xin nghỉ phép lên giám đốc Tanigo. Bất ngờ thay, vị giám đốc đáng kính không những hào phóng phê duyệt đơn nghỉ phép ngay lập tức, mà ông còn nảy ra ý định tuyệt vời là đài thọ cho toàn bộ nhân viên công ty cùng đi nghỉ mát chung để củng cố tinh thần đoàn kết.

``Quả đúng là một vị CEO có tâm nhất quả đất nhỉ?'' chắc hẳn các bạn sẽ tự hỏi như vậy. Và nếu nhìn từ góc độ phúc lợi nhân viên, thì sự thật hoàn toàn đúng là như thế!

\ct

Hai ngày sau, trên một bãi biển tuyệt đẹp ngập tràn ánh nắng vàng rực rỡ, tiếng sóng vỗ rì rào hòa quyện cùng tiếng cười đùa giòn giã của các cô gái đã tạo nên một bức tranh mùa hè sống động. Tất cả các thành viên \hll\ đã từng xuất hiện tại văn phòng trong suốt thời gian qua đều tề tựu đông đủ. Khoác lên mình những bộ đồ bơi rực rỡ và lộng lẫy nhất, họ cùng nhau tận hưởng từng khoảnh khắc tuyệt vời của kỳ nghỉ dưỡng.

Dưới bãi nước nông, Miko và Haato đang tươi cười tạo dáng, thực hiện những cảnh quay đầu tiên cho MV ca nhạc đặc biệt \emph{``Shiny Smily Story''}. Trên bờ cát, Kuon đang cầm cuốn kịch bản, chạy đôn chạy đáo hướng dẫn đội hình chuẩn bị cho phân cảnh vũ đạo chung của cả nhóm. Phía xa xa, Ayame và Matsuri đang vui vẻ té nước vào nhau, trong khi Aki và Fubuki hì hục chỉnh sửa lại đống đạo cụ quay phim.

Khi giai điệu sôi động của bài hát cất lên, mọi người cùng nắm tay nhau chạy ùa về phía biển, ánh mắt ai cũng ngập tràn sự tự do và hạnh phúc. Phân cảnh khép lại với hình ảnh toàn bộ đại gia đình \hll\ đứng xếp hàng ngang trên bờ cát trắng, quay mặt về phía ống kính và đồng thanh hô vang:

\begin{center}
    \uline{\textbf{``Hãy cùng nhau tạo nên những kỷ niệm tuyệt đẹp nhất trong mùa hè này nhé!''}}\\
    Các bạn có thể thưởng thức MV tại đây:
    \embedvideo{\includegraphics[width=12.5cm]{../../../../Image/Book1/1-10k.jpg}}{mv.mp4}

    \uline{Hoặc nhấn vào \href{https://www.youtube.com/watch?v=RrsGNMMghKM}{đường dẫn này} nếu video không hiển thị.}
\end{center}

\end{document}
